% =====================================================================
\section{Solutions to the Predicted Problems}\label{ch:predsol}
% =====================================================================

\subsection{Part A}

\solhead{Problem 1 (22 points).}
\textbf{(a)} $z=\frac43x$. Substituting:
$x^2+y^2+\frac{16}{9}x^2=100\iff\frac{25}{9}x^2+y^2=100\iff\frac{x^2}{36}+\frac{y^2}{100}=1$.
So $x=6\cos t$, $y=10\sin t$, $z=\frac43(6\cos t)=8\cos t$:
\[
  \boxed{\vr(t)=\ip{6\cos t,\;10\sin t,\;8\cos t},\quad0\le t\le2\pi.}
\]
Checks: $36\cos^2t+100\sin^2t+64\cos^2t=100$\checkok; $4(6\cos t)-3(8\cos t)=0$\checkok;
$\vr(0)=\ip{6,0,8}$\checkok; the shadow $(6\cos t,10\sin t)$ runs counterclockwise\checkok.
(Again a great circle of radius 10: $\vr=\cos t\ip{6,0,8}+\sin t\ip{0,10,0}$.)

\textbf{(b)} $R$: $6\cos t=3$, $10\sin t=5\sqrt3\Rightarrow t=\frac\pi3$ (and $8\cos\frac\pi3=4$\checkok).
$\vr\,'(t)=\ip{-6\sin t,\;10\cos t,\;-8\sin t}$, so
$\vr\,'\!\left(\frac\pi3\right)=\ip{-3\sqrt3,\;5,\;-4\sqrt3}$ and
\[
  \boxed{\mathbf L(s)=\ip{3,\,5\sqrt3,\,4}+s\ip{-3\sqrt3,\,5,\,-4\sqrt3}}.
\]

\textbf{(c)} $\norm{\vr\,'}=\sqrt{36\sin^2t+100\cos^2t+64\sin^2t}=10$. $S(-6,0,-8)$ is $t=\pi$. So
$L=\int_0^\pi10\,\dd t=\boxed{10\pi}$ (half of the great circle).

\textbf{(d)} $\vr\,'(0)=\ip{0,10,0}$, $\vr\,''(0)=\ip{-6,0,-8}$,
$\vr\,'\times\vr\,''=\ip{-80,0,60}$ with norm $100$; $\norm{\vr\,'}^3=1000$, so
$\kappa=\boxed{\frac{1}{10}}$ (circle of radius $10$).

\solhead{Problem 2 (15 points).}
Term by term: $\pd{}{x}$: $3y^2+\frac{2x}{x^2+y^2}+\frac{2y}{x^2}$ (since $-2yx^{-1}\to2yx^{-2}$);
$\pd{}{y}$: $6xy+\frac{2y}{x^2+y^2}-\frac2x$. At $(1,1)$: $f_x=3+1+2=6$, $f_y=6+1-2=5$, so
$\grad f(P)=\ip{6,5}$.

\textbf{(a)} $\vecPQ=\ip{3,-4}$, $\vu=\ip{\frac35,-\frac45}$: $\Du f=\frac{18-20}{5}=\boxed{-\frac25}$ (decreasing).

\textbf{(b)} $\boxed{\frac{1}{\sqrt{61}}\ip{6,5}}$.\quad
\textbf{(c)} $\boxed{\sqrt{61}}$.\quad
\textbf{(d)} $\vu\perp\ip{6,5}$: $\boxed{\pm\frac{1}{\sqrt{61}}\ip{5,-6}}$.

\solhead{Problem 3 (13 points).}
\textbf{(a)} $M_y=\cos(xy)-xy\sin(xy)=N_x$; $N_z=e^z=P_y$; $P_x=2x=M_z$. All continuous on
$\R^3$, so $\vF$ is conservative.

\textbf{(b)} $f=\int M\,\dd x=\sin(xy)+x^2z+g(y,z)$.
$f_y=x\cos(xy)+g_y=x\cos(xy)+e^z\Rightarrow g=ye^z+h(z)$.
$f_z=x^2+ye^z+h'(z)=x^2+ye^z\Rightarrow h=K$.
$f(0,0,0)=0+0+0+K=2$:
\[
  \boxed{f(x,y,z)=\sin(xy)+x^2z+ye^z+2}.
\]
\textbf{(c)} $\grad f(1,0,0)=\vF(1,0,0)=\ip{0+0,\;1+1,\;1+0}=\ip{0,2,1}$. Maximum rate
$\boxed{\sqrt5}$. Toward $(1,2,2)$: $\vu=\frac{1}{\sqrt2}\ip{0,1,1}$, $\Du f=\frac{2+1}{\sqrt2}=\boxed{\frac{3\sqrt2}{2}}$.

\solhead{Problem 4.}
$\vr(0)=\ip{1,0,1}$, so start at $t=0$.
$\vr\,'=e^t\ip{\cos t-\sin t,\;\sin t+\cos t,\;1}$ and
$\norm{\vr\,'}^2=e^{2t}\left[(\cos t-\sin t)^2+(\sin t+\cos t)^2+1\right]=e^{2t}(2+1)=3e^{2t}$, so
$\norm{\vr\,'}=\sqrt3\,e^t$ (not constant!).
$s(t)=\int_0^t\sqrt3e^u\,\dd u=\sqrt3(e^t-1)$. Setting $s=\sqrt3$: $e^t=2$, $t=\ln2$.
\[
  \vr(\ln2)=\boxed{\left(2\cos(\ln2),\;2\sin(\ln2),\;2\right)}\approx(1.538,\,1.278,\,2).
\]

\solhead{Problem 5.}
$\vr=\ip{t,\;b\sin t,\;0}$, $\vr\,'=\ip{1,b\cos t,0}$, $\vr\,''=\ip{0,-b\sin t,0}$. At $t=\frac\pi2$:
$\vr\,'=\ip{1,0,0}$, $\vr\,''=\ip{0,-b,0}$, $\vr\,'\times\vr\,''=\ip{0,0,-b}$. So
$\kappa=\frac{\abs b}{1}=2\Rightarrow\boxed{b=\pm2}$.

\solhead{Problem 6.}
$\grad f=\ip{2x+2y,\;2x}$, $\grad f(1,1)=\ip{4,2}$ ($\norm{\grad f}=2\sqrt5>2$, so two answers).
$4a+2b=2\Rightarrow b=1-2a$; $a^2+(1-2a)^2=1\iff5a^2-4a=0$:
$\boxed{\vu=\ip{0,1}\ \text{or}\ \vu=\ip{\frac45,-\frac35}}$. Check: $2$ and $\frac{16}{5}-\frac65=2$.\checkok

\solhead{Problem 7.}
$\grad f=\ip{2x-4,\;4y+4}=\lambda\ip{1,-1}$, $\lambda>0$: $2x-4=\lambda>0$ and $4y+4=-\lambda$.
So $4y+4=-(2x-4)\iff y=-\frac x2$, with $x>2$:
$\boxed{\left\{\left(x,-\tfrac x2\right):x>2\right\}}$ (the point $(2,-1)$, where $\grad f=\vzero$, is excluded).

\solhead{Problem 8.}
$M_y=2x$, $N_x=2bx\Rightarrow b=1$. $N_z=3$, $P_y=c\Rightarrow c=3$. $P_x=6z^2$, $M_z=3az^2\Rightarrow a=2$.
With $\vF=\ip{2xy+2z^3,\;x^2+3z,\;6xz^2+3y}$:
$f=x^2y+2xz^3+g(y,z)$; $f_y=x^2+g_y=x^2+3z\Rightarrow g=3yz+h(z)$;
$f_z=6xz^2+3y+h'=6xz^2+3y\Rightarrow h=K$.
$\boxed{a=2,\ b=1,\ c=3,\quad f=x^2y+2xz^3+3yz+K}$.

\subsection{Part B}

\solhead{Problem 9.}
\textbf{(a)} $x^2+\frac{y^2}{4}=1$: $\boxed{\vr(t)=\ip{\cos t,\;2\sin t,\;2}}$, $0\le t\le2\pi$.

\textbf{(b)} $\vr\,'=\ip{-\sin t,2\cos t,0}$, $\vr\,''=\ip{-\cos t,-2\sin t,0}$,
$\vr\,'\times\vr\,''=\ip{0,0,2\sin^2t+2\cos^2t}=\ip{0,0,2}$, and
$\norm{\vr\,'}^2=\sin^2t+4\cos^2t=4-3\sin^2t$. So
\[
  \kappa(t)=\frac{2}{(4-3\sin^2t)^{3/2}} .
\]
The major axis of this ellipse is along $y$ (semi-axis $2$), the minor axis along
$x$ (semi-axis $1$).
\begin{itemize}
\item Maximum when the denominator is smallest, $\sin^2t=1$: $\boxed{\kappa_{\max}=2}$ at
      $(0,\pm2,2)$, the ends of the \emph{major} axis, where the ellipse turns most sharply.
\item Minimum when $\sin t=0$: $\boxed{\kappa_{\min}=\frac{2}{4^{3/2}}=\frac14}$ at $(\pm1,0,2)$,
      the ends of the \emph{minor} axis, where it is flattest.
\end{itemize}
(Check with $\frac{a}{b^2}$ and $\frac{b}{a^2}$ for $a=1,b=2$: $\frac14$ and $2$.\checkok)

\textbf{(c)} $L=\displaystyle\int_0^{2\pi}\sqrt{\sin^2t+4\cos^2t}\,\dd t=\int_0^{2\pi}\sqrt{4-3\sin^2t}\,\dd t$
(an elliptic integral, $\approx9.69$).

\solhead{Problem 10.}
$\norm{\vr\,'}=\sqrt{4\sin^2t+4\cos^2t+5}=3$, $\vr(0)=(2,0,0)$, so $t=\frac{6\pi}{3}=2\pi$:
$\vr(2\pi)=\boxed{\left(2,\;0,\;2\sqrt5\,\pi\right)}$, exactly one turn of the helix higher.

\solhead{Problem 11.}
$\kappa(x)=\dfrac{e^x}{(1+e^{2x})^{3/2}}$. Differentiate:
\[
  \kappa'(x)=\frac{e^x(1+e^{2x})^{3/2}-e^x\cdot\frac32(1+e^{2x})^{1/2}\cdot2e^{2x}}{(1+e^{2x})^3}
  =\frac{e^x(1+e^{2x})^{1/2}\left(1-2e^{2x}\right)}{(1+e^{2x})^3}.
\]
$\kappa'=0\iff e^{2x}=\frac12\iff x=-\frac{\ln2}{2}$; $\kappa'>0$ before and $<0$ after, so it is a maximum.
Point: $\left(-\frac{\ln2}{2},\,\frac{1}{\sqrt2}\right)$; maximum curvature
$\kappa=\dfrac{1/\sqrt2}{(3/2)^{3/2}}=\boxed{\dfrac{2\sqrt3}{9}}\approx0.385$.

\solhead{Problem 12.}
$\grad f=\ip{e^{yz}+\frac{1}{x+y+z},\;xze^{yz}+\frac{1}{x+y+z},\;xye^{yz}+\frac{1}{x+y+z}}$, so
$\grad f(1,0,0)=\ip{2,1,1}$. Toward $(3,2,1)$: $\ip{2,2,1}/3$, $\Du f=\frac{4+2+1}{3}=\boxed{\frac73}$.
Maximum rate $\boxed{\sqrt6}$.

\solhead{Problem 13.}
$M_y=2y=N_x$ on the half-plane $y>0$ (simply connected): conservative.
$f=\ln(1+x^2)+xy^2+g(y)$; $f_y=2xy+g'(y)=2xy+\frac1y\Rightarrow g=\ln y$.
$\boxed{f=\ln(1+x^2)+xy^2+\ln y+K}$.

\solhead{Problem 14.}
$P_x=y+1$ but $M_z=y$; these differ everywhere, so $\vF$ is not conservative.
(The other two conditions hold: $M_y=z=N_x$, $N_z=x=P_y$, which is why all three must be checked.)

\solhead{Problem 15.}
\textbf{(a)} $\vF=\grad f=\ip{2xy,\;x^2+z^3,\;3yz^2}$. $M_y=2x=N_x$, $N_z=3z^2=P_y$, $P_x=0=M_z$.\checkok\\
\textbf{(b)} $\grad f(1,2,-1)=\ip{4,\;1-1,\;6}=\ip{4,0,6}$: direction $\frac{1}{\sqrt{13}}\ip{2,0,3}$, maximum
rate $\sqrt{52}=\boxed{2\sqrt{13}}$.\\
\textbf{(c)} $\vecPQ=\ip{1,2,2}$, length $3$: $\Du f=\frac{4+0+12}{3}=\boxed{\frac{16}{3}}$.

\solhead{Problem 16.}
$\vF(x,y)=\ip{y,x}$ vanishes only at the origin, has length $r$, is symmetric about
$y=x$, points away from the origin along $y=x$ (e.g. $\vF(1,1)=\ip{1,1}$) and toward it
along $y=-x$ ($\vF(1,-1)=\ip{-1,1}$): a saddle flow along the hyperbolas $x^2-y^2=c$.
It is a gradient field: $M_y=1=N_x$ and $\vF=\grad(xy)$.
